\note{sa page 1}

Soit $k$ anneau (commut.) fixé.

1. Soient \struck{$k$ un anneau,} $M^{*}$ un $k$-module gradué (à graduation
\ill{}), $\mathcal{E}^{*}$ un \add{sous-anneau de l'anneau} \struck{une algèbre
graduée} d'endomorphismes de $M^{*}$. On désigne par $\pi_{i}$ la projection
canonique de $M^{*}$ sur $M^{i}$. \struck{et} On suppose que l'on a :

\[ (1.\ill{}) \qquad \pi_{i} \in \mathcal{E}^{*} \quad \text{pour tout } i . \]

\note{L'étiquette de cette formule est réécrite par-dessus une autre et
n'est plus lisible ; le second chiffre pourrait être un $2$. Les étiquettes
(1.1), (1.2), (1.3), (1.4) sont, elles, sûres.}

Cela implique que $\mathcal{E}^{*}$ est un sous-\struck{algèbre}
\add{anneau} gradué\supplied{e} de l'algèbre graduée \struck{des}
$\operatorname{End}(M^{*})$ (dont la graduation est définie, si l'on veut, par
\[ \operatorname{End}^{(\nu)}(M) = \textstyle\sum \operatorname{Hom}(M^{i}, M^{i+\nu}) \]
ou de façon équivalente
\[ \operatorname{End}^{(\nu)}(M) = \textstyle\sum_{i} \pi_{i+\nu} \operatorname{End}(M^{*})\, \pi_{i} , \]
la composante homogène de degré $\nu$ d'un endomorphisme $u$ de $M^{*}$ étant
\[ \struck{\mathrm{gr}}\ \mathrm{pr}_{\nu}(u) = \textstyle\sum_{i} \pi_{i+\nu}\, u\, \pi_{i} . \]

La graduation de $\mathcal{E}^{*}$ est définie (intrinsèquement, en termes de la
donnée des $\pi_{i}$), par

\[ (1.3) \qquad \text{a)}\ \ \mathcal{E}^{(\nu)} = \sum_{i} \pi_{i+\nu}\, \mathcal{E}^{*}\, \pi_{i} , \qquad \text{b)}\ \ \mathrm{pr}_{\nu}\, u = \sum_{i} \pi_{i+\nu}\, u\, \pi_{i} , \]

et

\[ (1.4) \qquad u \in \mathcal{E}^{(\nu)} \iff u\, \pi_{i} = \pi_{i+\nu}\, u \quad \text{pour tout } i . \]

\marginal{Noter les relations}

\[ (1.1) \qquad \begin{cases} \pi_{i}^{2} = \pi_{i} \\ \pi_{i} \pi_{j} = 0 \ \text{ si } i \neq j \\ \sum_{i} \pi_{i} = 1 \end{cases} \]

2. Soit $\mathcal{E}^{*}$ un \struck{$k$-algèbre} anneau gradué, et soient
$\pi_{i}$ $(i \in \mathbb{Z})$ \add{nuls sauf un nombre fini} des éléments de
$\mathcal{E}$. À quelles conditions les formules (1.3) définissent-elles

